% Source: https://artofproblemsolving.com/community/c3h1641798_fun_problem
% Stable identity: aops-topic-1641798-q01
The fraction is defined only when $4x+6x^2\ne0$, so $x\ne0$ (and $x\ne-\dfrac{2}{3}$). Clearing the denominator,
\[
x^3y-2x^2+4x=4x+6x^2\quad\Longrightarrow\quad x^3y=8x^2\quad\Longrightarrow\quad xy=8 .
\]
Since $y>0$, also $x>0$, and $y=\dfrac{8}{x}$. The second equation becomes $\left(\dfrac{8}{x}\right)^x=16$, or, putting $x=\dfrac{8}{y}$, $y^{\dfrac{8}{y}}=16$. Taking logarithms,
\[
\dfrac{8\ln y}{y}=\ln16=4\ln2\qquad\Longleftrightarrow\qquad \dfrac{\ln y}{y}=\dfrac{\ln 2}{2}.
\]

\textbf{Two solutions.} $y=2$ works, and so does $y=4$, because $\dfrac{\ln4}{4}=\dfrac{2\ln2}{4}=\dfrac{\ln2}{2}$.

\textbf{No others.} Let $h(y)=\dfrac{\ln y}{y}$ for $y>0$. Then $h'(y)=\dfrac{1-\ln y}{y^2}$, which is positive for $y<e$ and negative for $y>e$. So $h$ is strictly increasing on $(0,e]$ and strictly decreasing on $[e,\infty)$, and it takes each value at most once on each of these intervals. Hence $y=2$ and $y=4$ are the only solutions.

They give $(x,y)=(4,2)$ and $(x,y)=(2,4)$; both check: $4\cdot2=2\cdot4=8$, $2^4=4^2=16$, and the denominators are nonzero. Therefore
\[
\dfrac{y}{x}=\boxed{\dfrac{1}{2}\ \text{or}\ 2}.
\]
